\chapter{Complexity — Run}
%%% generated by the blueprint pipeline from TCSlib.Complexity.SpaceComplexity.Codec.Run
\section{Overview}

Declarations of the Lean module \texttt{TCSlib.Complexity.SpaceComplexity.Codec.Run}.

\section{Declarations}

\begin{definition}[Two-language codec]
\lean{Complexity.Codec.codecAs}
\label{Complexity.Codec.codecAs}
\leanok
\uses{Complexity.inputBitLang, Complexity.inputLenLang}
Define a function from $	ext{Fin}(2)$ to Boolean languages by assigning the input-length
language to index $0$ and the input-bit language to index $1$.
\end{definition}

\begin{lemma}[Log-space decidability of the codec languages]
\lean{Complexity.Codec.codecAs_mem}
\label{Complexity.Codec.codecAs_mem}
\leanok
\uses{Complexity.Codec.codecAs, Complexity.LOGSPACE, Complexity.inputBitLang_mem, Complexity.inputLenLang_mem}
\difficulty{4}
For each of the two languages selected by the codec map—that is, the language of
encodings whose indexed position is within the first string's length, or the language of
encodings whose indexed bit is true—the language is decidable in logarithmic space.
\end{lemma}

\begin{definition}[Oracle for the code language]
\lean{Complexity.Codec.orc}
\label{Complexity.Codec.orc}
\leanok
\uses{Complexity.Codec.codecAs, Turing.MultiTapeTM.indicator}
For each tape index $j\in\{0,1\}$ and Boolean list $V$, the oracle returns the indicator
of whether $V$ belongs to the code set associated with $j$.
\end{definition}

\begin{lemma}[Bit oracle on encoded input]
\lean{Complexity.Codec.orc_bit}
\label{Complexity.Codec.orc_bit}
\leanok
\uses{Complexity.Codec.jBit, Complexity.Codec.orc, Complexity.indicator_inputBitLang, Complexity.inputBitLang, Turing.MultiTapeTM.indicator, Turing.pairEncode}
\difficulty{3}
For every bit string $y$ and natural number $a$, the bit oracle applied to the pair
encoding of $y$ and the binary bit string for $a$ returns the bit of $y$ at index $a$,
with false returned when $a$ is outside the string.
\end{lemma}

\begin{lemma}[Length oracle on encoded input]
\lean{Complexity.Codec.orc_len}
\label{Complexity.Codec.orc_len}
\leanok
\uses{Complexity.Codec.jLen, Complexity.Codec.orc, Complexity.indicator_inputLenLang, Complexity.inputLenLang, Turing.MultiTapeTM.indicator, Turing.pairEncode}
\difficulty{3}
For every bit string $y$ and natural number $a$, the length oracle applied to the pair
encoding of $y$ and the binary bit string for $a$ returns true exactly when $a<|y|$.
\end{lemma}

\begin{definition}[Bounded coded-machine reachability]
\lean{Complexity.Codec.Reach}
\label{Complexity.Codec.Reach}
\leanok
\uses{Complexity.Codec.CA, Complexity.Codec.CMode, Complexity.Codec.Lab, Complexity.Codec.orc, Complexity.LogProg.AReach, Complexity.LogProg.GuardLe, Turing.FinNDTM}
For a finite nondeterministic Turing machine $N$ over the Boolean alphabet, a mode $m$,
a Boolean list $y$, a natural-number bound $B$, labels $\ell,\ell'$, and vectors
$v,v'\colon\{0,\ldots,20\}\to\mathbb N$, the relation holds exactly when the compiled
machine for $N$ and $m$ can reach configuration $(\ell',v')$ from $(\ell,v)$ on input
$y$, using the specified code oracle and the guard determined by $B$.
\end{definition}

\begin{lemma}[Transitivity of bounded reachability]
\lean{Complexity.Codec.Reach.trans}
\label{Complexity.Codec.Reach.trans}
\leanok
\uses{Complexity.Codec.Lab, Complexity.Codec.Reach, Complexity.LogProg.AReach.trans}
\difficulty{1}
Fix a nondeterministic machine $N$, a mode, an input word $y$, and a bound $B$. If the
codec program has a run from configuration $(l_1,v_1)$ to $(l_2,v_2)$ and a run from
$(l_2,v_2)$ to $(l_3,v_3)$, with every configuration before the endpoint of each run
satisfying the program's guard and all 21 register values at most $B$, then it has such
a run from $(l_1,v_1)$ to $(l_3,v_3)$.
\end{lemma}

\begin{lemma}[Reachability under endpoint recasting]
\lean{Complexity.Codec.Reach.cast}
\label{Complexity.Codec.Reach.cast}
\leanok
\uses{Complexity.Codec.Lab, Complexity.Codec.Reach}
\difficulty{1}
Fix a finite nondeterministic Boolean machine, a computation mode, an input word, a
register bound $B$, two labels $l_1,l_2$, an initial register assignment $v_1$, and
final register assignments $v_2,v_2'$. If the compiled computation reaches label $l_2$
with registers $v_2$ from label $l_1$ with registers $v_1$, while satisfying the
reachability guard at every earlier step, and $v_2=v_2'$, then it also reaches label
$l_2$ with registers $v_2'$ under the same guard and from the same initial
configuration.
\end{lemma}

\begin{lemma}[Reachability from a guarded single step]
\lean{Complexity.Codec.Reach.single'}
\label{Complexity.Codec.Reach.single'}
\leanok
\uses{Complexity.Codec.CA, Complexity.Codec.Lab, Complexity.Codec.Reach, Complexity.Codec.orc, Complexity.LogProg.AReach.single, Complexity.LogProg.GuardLe, Complexity.LogProg.astep}
\difficulty{3}
Fix a finite nondeterministic Boolean machine, a computation mode, an input word $y$, a
register bound $B$, labels $l,l'$, and register assignments $v,v'$. If the configuration
at label $l$ with registers $v$ satisfies the guard—that is, its instruction meets the
required preconditions and every register is at most $B$—and one computation step takes
it to label $l'$ with registers $v'$, then that endpoint is reachable from the initial
configuration $(l,v)$ while the guard holds at every earlier step.
\end{lemma}

\begin{lemma}[Register bounds imply the guarded condition]
\lean{Complexity.Codec.gd_of}
\label{Complexity.Codec.gd_of}
\leanok
\uses{Complexity.Codec.CA, Complexity.Codec.Lab, Complexity.LogProg.GuardLe, Complexity.LogProg.PreS}
\difficulty{1}
Let $N$ be a finite nondeterministic Turing machine with Boolean symbols, let
$\mathrm{mode}$ be a codec mode, let $y$ be a Boolean word, let $B\in\mathbb{N}$, let
$l$ be a program label, and let $v$ assign a natural number to each of the 21 registers.
If every register value is at most $B$ and the configuration with label $l$, registers
$v$, and no pending answer satisfies the codec precondition on $y$, then that
configuration satisfies the guard with bound $B$.
\end{lemma}

\begin{lemma}[Distinct source and destination registers]
\lean{Complexity.Codec.cpDst_ne}
\label{Complexity.Codec.cpDst_ne}
\leanok
\uses{Complexity.Codec.CpId, Complexity.Codec.cpDst, Complexity.Codec.cpSrc}
\difficulty{4}
For the instruction labels
$A,B,Q1s,Cap,Q2,N1,Tn,E,Ts,F,Tw,U,Q1f,E2,Ts2,F2,Tw2,U2,Q1f2,Tpos,Q3,Q4,Q1a,Q2a,Uc,Tsc$,
respectively, the ordered pairs (destination register index, source register index) are
$(2,1),(3,1),(15,1),(19,0),(16,0),(5,2),(13,5),(7,2),(13,4),(8,7),(13,6),(14,8),(15,8),(9,3),(13,4),(10,9),(13,6),(14,10),(15,10),(13,11),(17,2),(18,3),(15,2),(16,2),(14,2),(13,4)$.
In every pair, the destination and source indices differ.
\end{lemma}

\begin{theorem}[Execution of a register-copy routine]
\lean{Complexity.Codec.cp_run}
\label{Complexity.Codec.cp_run}
\leanok
\uses{Complexity.Codec.CA, Complexity.Codec.CpId, Complexity.Codec.Reach, Complexity.Codec.cpDst, Complexity.Codec.cpDst_ne, Complexity.Codec.cpNext, Complexity.Codec.cpSrc, Complexity.Codec.orc, Complexity.LogProg.registerCopy_loop}
\difficulty{1}
Fix a finite nondeterministic Turing machine $N$ over Boolean symbols and a computation
mode, and let $y$ be an input word, $B$ a register bound, $c$ a copy instruction, and
$v$ an initial assignment of natural numbers to the $21$ registers. If every register
value in $v$ is at most $B$, and the destination register's initial value is at most the
source register's initial value, then the specified register program has a guarded
execution from the copy instruction and assignment $v$ to its prescribed continuation,
with the destination assigned the source's initial value and all other registers
unchanged.
\end{theorem}

\begin{lemma}[Distinct scan registers]
\lean{Complexity.Codec.scSrc_ne}
\label{Complexity.Codec.scSrc_ne}
\leanok
\uses{Complexity.Codec.ScId, Complexity.Codec.scDst, Complexity.Codec.scSrc}
\difficulty{4}
For each scan identifier $c\in\{E,F,E2,F2,In,C\}$, the source register, which is
register $13$, differs from the destination register: the destinations are registers
$7,8,9,10,14,14$, respectively.
\end{lemma}

\begin{theorem}[Scan execution and register update]
\lean{Complexity.Codec.sc_run}
\label{Complexity.Codec.sc_run}
\leanok
\uses{Complexity.Codec.CA, Complexity.Codec.Reach, Complexity.Codec.ScId, Complexity.Codec.orc, Complexity.Codec.scDst, Complexity.Codec.scFactor, Complexity.Codec.scIncs, Complexity.Codec.scNext, Complexity.Codec.scSrc, Complexity.Codec.scSrc_ne, Complexity.LogProg.registerScale_loop}
\difficulty{3}
Let $N$ be a finite nondeterministic Turing machine, choose a simulator mode and scan
identifier $c$, and let $y$ be a bit string with register bound $B$ and initial register
valuation $v$. If every register in $v$ is at most $B$, and the destination register's
initial value plus the scan factor for $c$ times the source register's initial value is
at most $B$, then the coded simulator has a finite guarded execution from scan state $c$
to its designated next state. At the end, the source register is zero, the destination
register has increased by the scan factor times the source register's initial value, and
all other registers retain their initial values; throughout the execution all registers
stay at most $B$ and the guard conditions hold.
\end{theorem}

\begin{theorem}[Increment-Count Register Run]
\lean{Complexity.Codec.ic_run}
\label{Complexity.Codec.ic_run}
\leanok
\uses{Complexity.Codec.CA, Complexity.Codec.IcId, Complexity.Codec.Reach, Complexity.Codec.icCount, Complexity.Codec.icL, Complexity.Codec.icNext, Complexity.Codec.icReg, Complexity.Codec.incChain, Complexity.Codec.orc}
\difficulty{3}
Let $N$ be a finite nondeterministic Turing machine with Boolean input alphabet, let $m$
be a codec mode, let $y$ be a bit string, let $B$ be a natural number, and let $c$ be an
increment-count identifier. For a register assignment $v$ on $21$ registers, suppose
every register value is at most $B$ and the value in the register associated with $c$,
plus the count associated with $c$, is at most $B$. Then the codec register program can
reach the next label for $c$ from its index-zero label, increasing that register by the
associated count; throughout the run, its prescribed safety condition holds and every
register remains at most $B$.
\end{theorem}

\begin{lemma}[Initial Increment Label]
\lean{Complexity.Codec.icL_zero}
\label{Complexity.Codec.icL_zero}
\leanok
\uses{Complexity.Codec.IcId, Complexity.Codec.W, Complexity.Codec.icCount, Complexity.Codec.icCount_le, Complexity.Codec.icL}
\difficulty{4}
Let $N$ be a finite nondeterministic Turing machine with Boolean input alphabet, and let
$c$ be an increment-count identifier. The label assigned to $c$ at index $0$ is its
indexed increment label with index $0$.
\end{lemma}

\begin{lemma}[Distinct decode registers]
\lean{Complexity.Codec.dc_distinct}
\label{Complexity.Codec.dc_distinct}
\leanok
\uses{Complexity.Codec.DcId, Complexity.Codec.dcAcc, Complexity.Codec.dcBase, Complexity.Codec.dcPtr, Complexity.Codec.rT}
\difficulty{4}
For each decode identifier $c\in\{S,\mathrm{Pos},\mathrm{Pos2}\}$, the pointer, base,
accumulator, and temporary registers are pairwise distinct. Explicitly, the pointer
differs from the base, accumulator, and temporary register, the base differs from the
accumulator and temporary register, and the accumulator differs from the temporary
register.
\end{lemma}

\begin{theorem}[Decode loop execution]
\lean{Complexity.Codec.dc_run}
\label{Complexity.Codec.dc_run}
\leanok
\uses{Complexity.Codec.CA, Complexity.Codec.DcId, Complexity.Codec.Reach, Complexity.Codec.Reach.cast, Complexity.Codec.dcAcc, Complexity.Codec.dcBase, Complexity.Codec.dcDowns, Complexity.Codec.dcIncs, Complexity.Codec.dcNext, Complexity.Codec.dcPtr, Complexity.Codec.dcStride, Complexity.Codec.dc_distinct, Complexity.Codec.decodeStride_loop, Complexity.Codec.jBit, Complexity.Codec.orc, Complexity.Codec.orc_bit, Complexity.Codec.rT, Complexity.LogProg.bitsVal}
\difficulty{5}
Fix a finite nondeterministic Turing machine $N$, a mode, a bit string $y$, a bound $B$,
a decode identifier $c$, a bit string $w$, and a register assignment $v$ to the 21
registers. Suppose every suffix value of $w$ is at most $B$; every register value under
$v$ is at most $B$; the pointer register equals the base register plus the stride for
$c$ times $|w|$; the accumulator and temporary registers are zero; and, at each position
$i<|w|$, the input bit at the base offset by the stride times $i$ equals the $i$th bit
of $w$. Then the codec program in the given mode, starting at its decode-test
instruction with assignment $v$, reaches its next instruction for $c$ with the pointer
reset to its base value and the accumulator equal to the binary value of $w$, leaving
all other registers unchanged. Every configuration before the endpoint satisfies the
program's safety condition and has all register values at most $B$.
\end{theorem}

\begin{theorem}[Input length scan]
\lean{Complexity.Codec.len_run}
\label{Complexity.Codec.len_run}
\leanok
\uses{Complexity.Codec.CA, Complexity.Codec.Reach, Complexity.Codec.jLen, Complexity.Codec.orc, Complexity.Codec.orc_len, Complexity.Codec.rL, Complexity.LogProg.inputLength_loop}
\difficulty{1}
Let $N$ be a finite-state nondeterministic Turing machine with Boolean input, let $y$ be
a Boolean word, and let $B$ be a natural number with $|y|\le B$. Suppose the initial
register contents $v$ satisfy $v(r)\le B$ for every register and $v(r_L)\le |y|$. In the
register program with mode $\mathrm{mode}$, there is a run from the length-scan
instruction to the first prefix-scan instruction that leaves every register unchanged
except $r_L$, which is set to $|y|$; every configuration before the final one satisfies
the run guard with bound $B$.
\end{theorem}

\begin{theorem}[Prefix scan of a paired word]
\lean{Complexity.Codec.pr_run}
\label{Complexity.Codec.pr_run}
\leanok
\uses{Complexity.Codec.CA, Complexity.Codec.Reach, Complexity.Codec.jBit, Complexity.Codec.orc, Complexity.Codec.orc_bit, Complexity.Codec.prAdvL, Complexity.Codec.prCont, Complexity.Codec.rL, Complexity.Codec.rP, Complexity.LogProg.alignedPair_loop, Complexity.getD_drop, Turing.pairDecode}
\difficulty{4}
Let $N$ be a finite nondeterministic Turing machine over Boolean symbols, let $mode$ be
one of the four codec modes, let $y$ be a Boolean word, and let $B$ be a natural-number
register bound. Let $p$ be one of three scan indices and let $v$ be a register
assignment such that every register is at most $B$, the length register equals $|y|$,
and the position register is at most $|y|$. If decoding the suffix of $y$ beginning at
the position register fails, then for some $t$ no greater than the length of that
suffix, the scan reaches rejection after increasing the position register by $t$, with
all other registers unchanged. If decoding succeeds with first component $u$, the scan
reaches its prescribed continuation after increasing the position register by $2|u|+2$,
leaving all other registers unchanged.
\end{theorem}

\begin{theorem}[Execution of the capacity-check routine]
\lean{Complexity.Codec.ct_run}
\label{Complexity.Codec.ct_run}
\leanok
\uses{Complexity.Codec.CA, Complexity.Codec.Reach, Complexity.Codec.jBit, Complexity.Codec.orc, Complexity.Codec.orc_bit, Complexity.Codec.rCap, Complexity.Codec.rL, Complexity.Codec.rQ1, Complexity.LogProg.canonicalTail_loop}
\difficulty{1}
Fix a finite nondeterministic Turing machine $N$ over Boolean symbols and a computation
mode. Let $y$ be an input word, $B$ a register bound, $w$ a Boolean word, and $v$ an
initial assignment of natural numbers to the $21$ registers. Suppose every register
value is at most $B$, the length register equals the position register plus the length
of $w$, and, at each position $i<|w|$, the input bit at the position register plus $i$
equals the $i$th bit of $w$ (with false used for out-of-range reads). Then there is an
assignment $v'$ such that every register remains at most $B$, the position register does
not increase, and every register other than the position and capacity registers is
unchanged. The specified register program has a guarded execution from the
capacity-check routine and $v$ to the state determined by whether $|w|$ is at most the
binary length of the capacity register's initial value and whether the last bit of $w$
(defaulting to true when $w$ is empty) is true: it continues to the designated next
routine exactly when both conditions hold, and otherwise rejects.
\end{theorem}

\begin{theorem}[Halving and Bit-Length Accumulation]
\lean{Complexity.Codec.hc_run}
\label{Complexity.Codec.hc_run}
\leanok
\uses{Complexity.Codec.CA, Complexity.Codec.Reach, Complexity.Codec.halvCount_loop, Complexity.Codec.orc, Complexity.Codec.rT, Complexity.Codec.rWp}
\difficulty{1}
Let $N$ be a finite nondeterministic Turing machine with Boolean input alphabet, let $m$
be one of the four codec modes, and let $y$ be a bit string and $B$ a natural number.
For any assignment $v$ of natural numbers to the $21$ registers, suppose every register
value is at most $B$ and the value in register $r_{
m Wp}$ plus the length of the binary
expansion of the value in register $r_{
m T}$ is at most $B$. Then the codec register
program can reach its copy phase for $E$ from its halving-counter test state, changing
register $r_{
m T}$ to $0$ and register $r_{
m Wp}$ to its former value plus that binary
expansion length; throughout the run, its prescribed safety condition holds and every
register remains at most $B$.
\end{theorem}

\begin{theorem}[Register comparison run]
\lean{Complexity.Codec.le_run}
\label{Complexity.Codec.le_run}
\leanok
\uses{Complexity.Codec.CA, Complexity.Codec.Reach, Complexity.Codec.le_loop, Complexity.Codec.orc, Complexity.Codec.rN1, Complexity.Codec.rPos, Complexity.Codec.rT}
\difficulty{3}
Let $N$ be a finite-state nondeterministic Turing machine with Boolean input, let $y$ be
a Boolean word, and let $B$ be a natural-number register bound. Suppose the initial
register contents $v$ satisfy $v(r)\le B$ for every register $r$ and $v(r_T)=0$. In the
register program with mode $\mathrm{mode}$, starting at the comparison instruction,
there is a run to the comparison’s next instruction: it reaches the clearing instruction
if $v(r_{\mathrm{Pos}})\le v(r_{\mathrm{N1}})$ and the rejecting instruction otherwise,
and the resulting registers agree with $v$ except that register $r_T$ contains
$\min(v(r_{\mathrm{Pos}}),v(r_{\mathrm{N1}}))$; every configuration before the final one
satisfies the run guard with bound $B$.
\end{theorem}

\begin{theorem}[Bounded Run of the Verifier]
\lean{Complexity.Codec.v_run}
\label{Complexity.Codec.v_run}
\leanok
\uses{Complexity.Codec.CA, Complexity.Codec.Reach, Complexity.Codec.VAd, Complexity.Codec.jBit, Complexity.Codec.orc, Complexity.Codec.orc_bit, Complexity.Codec.rE, Complexity.Codec.rQ1, Complexity.Codec.rQ2, Complexity.Codec.scan_loop, Complexity.Codec.vA1, Complexity.Codec.vA2, Complexity.Codec.δV}
\difficulty{4}
Let $N$ be a finite nondeterministic Turing machine with Boolean tape symbols, let the
codec mode be fixed, and let $y,u,u'$ be finite Boolean words and $B\in\mathbb N$.
Suppose $|u'|=|u|$, the starting verifier accumulator is $st$, and the starting register
values $v$ satisfy $v(r)\le B$ for every register $r$, $v(E)=v(Q1)+2|u|$, and
$v(Q2)+2|u|\le B$. Suppose also that for every $i<|u|$, the bits of $y$ at positions
$v(Q1)+2i$ and $v(Q2)+2i$ are respectively the $i$th bits of $u$ and $u'$, with
out-of-range bits interpreted as false. Then the codec program has a finite run, staying
at every time strictly before its endpoint in configurations satisfying its precondition
and register bound $B$, from label $vTest(st)$ with registers $v$ to label $vDone$ with
verifier state obtained by processing the zipped pairs of $u$ and $u'$ from $st$; the
final registers are $v$ with $Q1$ and $Q2$ updated respectively to $v(Q1)+2|u|$ and
$v(Q2)+2|u|$.
\end{theorem}

\begin{theorem}[Paired scan execution]
\lean{Complexity.Codec.s_run}
\label{Complexity.Codec.s_run}
\leanok
\uses{Complexity.Codec.CA, Complexity.Codec.Reach, Complexity.Codec.SAd, Complexity.Codec.SInfo, Complexity.Codec.jBit, Complexity.Codec.orc, Complexity.Codec.orc_bit, Complexity.Codec.rE, Complexity.Codec.rQ3, Complexity.Codec.rQ4, Complexity.Codec.sA1, Complexity.Codec.sA2, Complexity.Codec.scan_loop, Complexity.Codec.δS}
\difficulty{4}
Let $N$ be a finite nondeterministic Turing machine, let $y$ be a bit string, let $B$ be
a natural-number register bound, and let $I$ be a simulator state with two Boolean flags
$f_0,f_1$. Suppose $u,u'$ are bit strings of equal length, the register valuation $v$
has every register at most $B$, and its registers $E,Q_3,Q_4$ satisfy $v(E)=v(Q_3)+2|u|$
and $v(Q_4)+2|u|\le B$. Suppose also that, for every $i<|u|$, the input bits at
positions $v(Q_3)+2i$ and $v(Q_4)+2i$ in $y$ (taken to be false when out of range) equal
the $i$th bits of $u$ and $u'$, respectively (also taken false when out of range). Then
the coded simulator has a finite execution from the paired scan-test state, with
valuation $v$, to the paired scan-finish state whose simulator state is obtained by
folding the paired transition rule for $I$ over the zipped lists $u,u'$, and whose
valuation agrees with $v$ except that registers $Q_3,Q_4$ are advanced to $v(Q_3)+2|u|$
and $v(Q_4)+2|u|$. Throughout the execution, the simulator's guard conditions hold and
all registers remain at most $B$.
\end{theorem}
